\section{Estimates for class M1}\label{app:short-run-estimates}

This appendix proves the angle gap, \cref{prop:m1-angle-gap}, and the arithmetic of the two-centre
case. Throughout, $K_*=97969/529200$ is the competitive cap of \cref{cor:casimir-cap},
$U(s)=1-s+s^2$, and $\omega(x)=-x-\log(1-x)$ for $0\le x<1$.

\subsection{Logarithm bounds}

\begin{lemma}\label{lem:m1-log}
  \begin{enumerate}[label=\textup{(\arabic*)}]
    \item If $z,z_0>0$ and $\log z_0\le\ell$, then $\log z\le\ell+z/z_0-1$.
    \item For $0<z\le1$ and $\tau=(1-z)/(1+z)$, $\log z\le-2\tau-\frac23\tau^3$.
    \item For $0<x\le1$, $\frac12(x-x^{-1})\le\log x$.
    \item If $a,b,c\ge0$, $c<1$ and $a^2+b^2\le c^2$, then $\omega(a)+\omega(b)\le\omega(c)$.
  \end{enumerate}
\end{lemma}

\begin{proof}
  (1) is $\log(z/z_0)\le z/z_0-1$. (2) keeps the first two terms of the series
  $\log z=-2\sum_{n\ge0}\tau^{2n+1}/(2n+1)$, all of whose terms are nonpositive. (3) is
  $y\le\sinh y$ for $y=-\log x\ge0$. (4) Write $\omega(x)=\sum_{n\ge0}x^{n+2}/(n+2)$. Since
  $a,b\le c$, termwise $a^{n+2}+b^{n+2}\le(a^2+b^2)c^n\le c^{n+2}$.
\end{proof}

Bound (2) gives the three rational values used below,
\begin{equation}\label{eq:m1-logs}
  \log\tfrac{17}{25}\le-\tfrac{10712}{27783},\qquad
  \log\tfrac34\le-\tfrac{296}{1029},\qquad
  \log\tfrac{17}{20}\le-\tfrac{8232}{50653},
\end{equation}
for $\tau=4/21$, $1/7$ and $3/37$.

\subsection{One reverse switch}

The reverse-switch term $\Psi(r,u)=r-\log r-(1-u)/r$ of the stationarity identity increases with
$u$ and, at fixed $u$, is largest when the switch bound $u\le U(r)$ is attained.

\begin{lemma}\label{lem:m1-one-switch}
  Let $0<r<1/2$ and $3/4\le u\le U(r)$, and put $s=\sqrt{4u-3}$. Then $0\le s<1$, $2r\le 1-s$, and
  \[
    r-\log r-\frac{1-u}r\le\omega(s)+\log2 .
  \]
\end{lemma}

\begin{proof}
  Put $r'=(1-s)/2$, so that $U(r')=(3+s^2)/4=u$. Since $U$ decreases on $[0,1/2]$ and
  $U(r')=u\le U(r)$, we have $r\le r'$. The identity
  \[
    2r'-1-\Bigl(r-\tfrac12\bigl(\tfrac r{r'}-\tfrac{r'}r\bigr)-\frac{1-U(r')}r\Bigr)
    =\frac{(r'-r)^2(1-2r')}{2rr'}
  \]
  and $1-2r'=s\ge0$ give $r-\frac12(r/r'-r'/r)-(1-u)/r\le 2r'-1=-s$. By
  \cref{lem:m1-log}(3) with $x=r/r'\le1$, $-\log(r/r')\le-\frac12(r/r'-r'/r)$, so
  $r-\log r-(1-u)/r\le -s-\log r'$. Finally $-\log r'=-\log(1-s)+\log2$.
\end{proof}

\subsection{Reduction to two variables}

\begin{lemma}\label{lem:m1-reduce}
  Let $\gamma^2,r,t,\rho,u,v,E$ satisfy \textup{(i)--(iv)} of \cref{lem:m1-centre-data}, and suppose
  that $E\ge\frac35\sqrt{uv}$. Let $s=\sqrt{4u-3}$, $s'=\sqrt{4v-3}$,
  $S=\sqrt{s^2+s'^2+s^2s'^2/3}$ and $z=4\rho/(1-S)$. Then $0\le S<1$, $0<z\le1$,
  \[
    uv=\frac{3(3+S^2)}{16},\qquad \rho=\frac{(1-S)z}4,
  \]
  \[
    E\le(1-S)\Bigl(1-\frac z4\Bigr)+\log z,\qquad E^2\ge\frac9{25}\cdot\frac{3(3+S^2)}{16}.
  \]
\end{lemma}

\begin{proof}
  The identity \eqref{eq:m1-stat} splits as
  $E=1-\rho+\log\rho+\Psi(r,u)+\Psi(t,v)$, and \cref{lem:m1-one-switch} bounds the two last terms
  by $\omega(s)+\log2$ and $\omega(s')+\log2$, with $2r\le1-s$ and $2t\le1-s'$. Then
  $4\rho<4rt\le(1-s)(1-s')$. Since $u,v\ge3/4$, $\sqrt{uv}\ge3/4$ and $E\ge9/20$, and the bounds
  just obtained give $9/20\le E<1-\rho-s-s'$, so $s+s'<11/20$. From $u=(3+s^2)/4$ and
  $v=(3+s'^2)/4$ one gets $uv=3(3+S^2)/16$. Next,
  \[
    (s+s'-ss')^2-S^2=\tfrac23s^2s'^2+2ss'(1-s-s')\ge0,
  \]
  so $S\le s+s'-ss'$ and $(1-s)(1-s')\le1-S$; in particular $S<1$ and $4\rho<1-S$, so
  $0<z<1$. By \cref{lem:m1-log}(4), $\omega(s)+\omega(s')\le\omega(S)$, and therefore
  \[
    E\le1-\rho+\log(4\rho)+\omega(S)=1-S-\rho+\log\frac{4\rho}{1-S}
    =(1-S)\Bigl(1-\frac z4\Bigr)+\log z .
  \]
  The last inequality of the lemma is the square of $E\ge\frac35\sqrt{uv}$.
\end{proof}

\begin{lemma}\label{lem:m1-core}
  Let $S\ge0$, $0<z\le1$ and $E>0$ satisfy
  $E\le(1-S)(1-z/4)+\log z$ and $E^2\ge\frac9{25}Q$, where $Q=3(3+S^2)/16$. Then
  \[
    16K_*\Bigl(1-\frac{(1-S)z}4\Bigr)^2<4Q-E^2 .
  \]
\end{lemma}

\begin{proof}
  The hypotheses give $E\ge\frac35\cdot\frac34=\frac9{20}$.

  \emph{The domain.} By $\log z\le z-1$, $E\le-S+z(3+S)/4\le\frac34(1-S)$, so $S\le2/5$. If
  $z<31/50$, \cref{lem:m1-log}(1) at $z_0=17/25$ with \eqref{eq:m1-logs} gives
  $E\le-\frac{10712}{27783}+z\bigl(\frac{25}{17}-\frac14\bigr)<\frac9{20}$, a contradiction. So
  $z\in[31/50,1]$ and $S\in[0,2/5]$.

  \emph{A supporting function on each interval.} Cover $[31/50,1]$ by the four intervals of
  \cref{tab:m1-core}. On an interval $[z_-,z_+]$ with tangent point $z_k$, bound $\ell_k$ and
  multiplier $\mu=\mu(S)\ge0$ from the table, let
  $T(z)=(1-S)(1-z/4)+\ell_k+z/z_k-1$, so that $0<E\le T(z)$ by \cref{lem:m1-log}(1). Since
  $E^2-\frac9{25}Q\ge0$,
  \begin{align*}
    4Q-E^2&\ge 4Q-E^2-\mu\bigl(E^2-\tfrac9{25}Q\bigr)
    \ge\bigl(4+\tfrac9{25}\mu\bigr)Q-(1+\mu)T(z)^2\\
    &=G(z)+16K_*\Bigl(1-\frac{(1-S)z}4\Bigr)^2,
  \end{align*}
  where the last equality defines $G$. It suffices to show $G(z)>0$ on the interval.

  \emph{Concavity.} $G$ is a constant minus two squares of affine functions of $z$, and
  \begin{multline*}
    (z_+-z_-)G(z)=(z_+-z)G(z_-)+(z-z_-)G(z_+)\\
    +(z_+-z_-)\Bigl[(1+\mu)\Bigl(\frac1{z_k}-\frac{1-S}4\Bigr)^2+K_*(1-S)^2\Bigr](z-z_-)(z_+-z).
  \end{multline*}
  So $G>0$ on $[z_-,z_+]$ once $G(z_-)>0$ and $G(z_+)>0$.

  \emph{The eight endpoint values.} Each $G(z_\pm)$ is a polynomial
  $c_0+c_1S+c_2S^2+c_3S^3$ in $S\in[0,2/5]$ with rational coefficients and $c_3\le0$. On
  $[0,2/5]$, $c_3S^3\ge\frac25c_3S^2$, so $G(z_\pm)\ge c_0+c_1S+c_2'S^2$ with
  $c_2'=c_2+\frac25c_3$. Replacing $c_0,c_1,c_2'$ by the smaller numbers $\underline c_0$,
  $\underline c_1$, $\underline c_2$ of \cref{tab:m1-core-ends} only lowers this quadratic for
  $S\ge0$, and each of the eight quadratics $\underline c_0+\underline c_1S+\underline c_2S^2$ is
  positive on $[0,2/5]$ by the reason given in the table: all three coefficients are nonnegative
  and $\underline c_0>0$; or $\underline c_2>0$ and the discriminant
  $\underline c_1^2-4\underline c_0\underline c_2$ is negative; or $\underline c_1\ge0$,
  $\underline c_2<0$ and the value $\underline c_0+\frac4{25}\underline c_2$ at $S=2/5$ is
  positive.
\end{proof}

\begin{table}[htbp]
  \centering
  \caption{The tangent lines and multipliers of \cref{lem:m1-core}.}\label{tab:m1-core}
  \begin{tabular}{@{}cccc@{}}
    \toprule
    interval $[z_-,z_+]$ & tangent point $z_k$ & bound $\ell_k\ge\log z_k$ & multiplier $\mu(S)$\\
    \midrule
    $[31/50,\,7/10]$ & $17/25$ & $-10712/27783$ & $41/200+9S/20$\\
    $[7/10,\,4/5]$ & $3/4$ & $-296/1029$ & $3/50+3S/5$\\
    $[4/5,\,9/10]$ & $17/20$ & $-8232/50653$ & $1/100+3S/20$\\
    $[9/10,\,1]$ & $1$ & $0$ & $0$\\
    \bottomrule
  \end{tabular}
\end{table}

\begin{table}[htbp]
  \centering
  \caption{Lower bounds for the coefficients of the quadratic minorants of $G(z_\pm)$ in the
  proof of \cref{lem:m1-core}, and the reason why each minorant is positive on
  $[0,2/5]$.}\label{tab:m1-core-ends}
  \begin{tabular}{@{}cccccc@{}}
    \toprule
    interval & end & $\underline c_0$ & $\underline c_1$ & $\underline c_2$ & positivity\\
    \midrule
    $[31/50,\,7/10]$ & $31/50$ & $0.0105$ & $0.0091$ & $-0.0018$ & value at $2/5$\\
                     & $7/10$  & $0.0105$ & $0.0691$ & $0.0907$ & coefficients\\
    $[7/10,\,4/5]$   & $7/10$  & $0.0112$ & $-0.0435$ & $0.2606$ & discriminant\\
                     & $4/5$   & $0.0110$ & $-0.0456$ & $0.3756$ & discriminant\\
    $[4/5,\,9/10]$   & $4/5$   & $0.0181$ & $-0.0326$ & $0.0903$ & discriminant\\
                     & $9/10$  & $0.0178$ & $-0.0193$ & $0.1181$ & discriminant\\
    $[9/10,\,1]$     & $9/10$  & $0.0153$ & $0.0132$ & $-0.0006$ & value at $2/5$\\
                     & $1$     & $0.0213$ & $0.0142$ & $0.0023$ & coefficients\\
    \bottomrule
  \end{tabular}
\end{table}

\begin{proof}[Proof of \cref{prop:m1-angle-gap}]
  Suppose $E\ge\frac35\sqrt{uv}$. \Cref{lem:m1-reduce} puts $(S,z,E)$ in the hypotheses of
  \cref{lem:m1-core}, whose conclusion, with $Q=uv$ and $(1-S)z/4=\rho$, reads
  $16K_*(1-\rho)^2<4uv-E^2=16(1-\rho)^2\gamma^2$. This contradicts $\gamma^2\le K_*$.
\end{proof}

The margin is small: on the region of \cref{fig:short-runs-gap} the least Casimir is about
$0.1863$, against $K_*\approx0.18513$. This is why the logarithm is replaced by four tangent lines
and not by one.

\subsection{The two-centre arithmetic}

\begin{lemma}\label{lem:m1-bc}
  If $0<a<b$, $a<c$ and $(b-a)(c-a)=2a^2$, then $bc>5.828\,a^2$ and $(a+b)(a+c)=2bc$.
\end{lemma}

\begin{proof}
  With $x=b/a$ and $y=c/a$, $(x-1)(y-1)=2$, so $xy=x+y+1$ and, by the inequality between the
  arithmetic and geometric means, $x+y-2\ge2\sqrt{(x-1)(y-1)}=2\sqrt2$. Hence
  $xy\ge3+2\sqrt2=5.8284\ldots>5.828$. Expanding, $(a+b)(a+c)=a^2+a(b+c)+bc$, and
  $bc-a(b+c)-a^2=(b-a)(c-a)-2a^2=0$.
\end{proof}
